% GENERATED FILE. Edit canonical lessons, then run npm run generate:books.
% canonical-source-hash: fnv1a64:ece62abc18bb6cfa
% canonical-lessons: IT-C315-celebrare, IT-C315-ricorrenza, IT-C315-anniversario, IT-C315-auguri, IT-C315-augurare

\chapter{Tanti auguri}
\label{ch:it-tanti-auguri}
\hlchaptermodality{315}

\begin{chapteropening}
\textbf{By the end of this chapter:} \emph{I can mark an occasion: celebrating, an anniversary, best wishes and wishing someone well.}

Complete the last lesson of chapter 315: I can mark an occasion: celebrating, an anniversary, best wishes and wishing someone well.
\end{chapteropening}

\section[celebrare]{celebrare — to celebrate formally; to hold (a ceremony)}

\label{lesson:IT-C315-celebrare}



\begin{quote}
Before the new one: say the Italian for the heir, then the Italian for the grave.
\end{quote}

\subsection*{You'll want to know: celebrare}
\textbf{celebrare} — \textquotedblleft{}to celebrate formally; to hold (a ceremony)\textquotedblright{}.

\begin{quote}
Abbiamo celebrato i cento anni della scuola.
\end{quote}

\begin{quote}
\emph{We celebrated the school's hundredth anniversary.}
\end{quote}

More formal than \emph{festeggiare}.

\begin{cousinweb}
\emph{celebrare} is a learned borrowing from Latin \emph{celebrāre} \textquotedblleft{}to frequent; to honour with a crowd\textquotedblright{} (Lewis \& Short), from \emph{celeber} \textquotedblleft{}crowded\textquotedblright{}.
\end{cousinweb}

\begin{grammarlens}[title={What you've built}]
One more word for telling, explaining and describing at B1.
\end{grammarlens}

\subsection*{Your turn}
\begin{itemize}
\raggedright
  \item \emph{Say it:} \emph{celebrare}
  \item \emph{Say it:} \emph{celebrare}, once more
  \item \emph{Say it:} a sentence of your own with \emph{celebrare}
  \item \emph{Recall:} say the Italian for the heir, then the Italian for the grave, then say \emph{celebrare} again
\end{itemize}

\begin{tcolorbox}[breakable,colback=teal!4,colframe=teal!35!black,title={Before you move on}]
What does celebrare mean? (\textquotedblleft{}To celebrate formally; to hold (a ceremony)\textquotedblright{}.) Say it once more.
\end{tcolorbox}
\section[la ricorrenza]{la ricorrenza — the recurring occasion, the anniversary date}

\label{lesson:IT-C315-ricorrenza}



\begin{quote}
Before the new one: say the Italian for the grave, then the Italian for to celebrate formally.
\end{quote}

\subsection*{You'll want to know: la ricorrenza}
\textbf{la ricorrenza} — \textquotedblleft{}the recurring occasion, the anniversary date\textquotedblright{}.

\begin{quote}
Il compleanno è una ricorrenza importante.
\end{quote}

\begin{quote}
\emph{A birthday is an important occasion.}
\end{quote}

Built on \emph{ricorrere}, ‘to come round again’.

\begin{cousinweb}
\emph{ricorrenza} is built on \emph{ricorrere} \textquotedblleft{}to recur\textquotedblright{}, from Latin \emph{recurrere} \textquotedblleft{}to run back\textquotedblright{} (Lewis \& Short).
\end{cousinweb}

\begin{grammarlens}[title={What you've built}]
One more word for telling, explaining and describing at B1.
\end{grammarlens}

\subsection*{Your turn}
\begin{itemize}
\raggedright
  \item \emph{Say it:} \emph{la ricorrenza}
  \item \emph{Say it:} \emph{la ricorrenza}, once more
  \item \emph{Say it:} a sentence of your own with \emph{la ricorrenza}
  \item \emph{Recall:} say the Italian for the grave, then the Italian for to celebrate formally, then say \emph{la ricorrenza} again
\end{itemize}

\begin{tcolorbox}[breakable,colback=teal!4,colframe=teal!35!black,title={Before you move on}]
What does la ricorrenza mean? (\textquotedblleft{}The recurring occasion, the anniversary date\textquotedblright{}.) Say it once more.
\end{tcolorbox}
\section[l'anniversario]{l'anniversario — the anniversary}

\label{lesson:IT-C315-anniversario}



\begin{quote}
Before the new one: say the Italian for to celebrate formally, then the Italian for the recurring occasion.
\end{quote}

\subsection*{You'll want to know: l'anniversario}
\textbf{l'anniversario} — \textquotedblleft{}the anniversary\textquotedblright{}.

\begin{quote}
Oggi è il nostro anniversario di matrimonio.
\end{quote}

\begin{quote}
\emph{Today is our wedding anniversary.}
\end{quote}

\emph{Buon anniversario!}

\begin{cousinweb}
\emph{anniversario} is a learned borrowing from Latin \emph{anniversārius} \textquotedblleft{}returning yearly\textquotedblright{} (EtymDB), from \emph{annus} \textquotedblleft{}year\textquotedblright{} and \emph{vertere} \textquotedblleft{}to turn\textquotedblright{}.
\end{cousinweb}

\begin{grammarlens}[title={What you've built}]
One more word for telling, explaining and describing at B1.
\end{grammarlens}

\subsection*{Your turn}
\begin{itemize}
\raggedright
  \item \emph{Say it:} \emph{l'anniversario}
  \item \emph{Say it:} \emph{l'anniversario}, once more
  \item \emph{Say it:} a sentence of your own with \emph{l'anniversario}
  \item \emph{Recall:} say the Italian for to celebrate formally, then the Italian for the recurring occasion, then say \emph{l'anniversario} again
\end{itemize}

\begin{tcolorbox}[breakable,colback=teal!4,colframe=teal!35!black,title={Before you move on}]
What does l'anniversario mean? (\textquotedblleft{}The anniversary\textquotedblright{}.) Say it once more.
\end{tcolorbox}
\section[gli auguri]{gli auguri — best wishes, greetings}

\label{lesson:IT-C315-auguri}



\begin{quote}
Before the new one: say the Italian for the recurring occasion, then the Italian for the anniversary.
\end{quote}

\subsection*{You'll want to know: gli auguri}
\textbf{gli auguri} — \textquotedblleft{}best wishes, greetings\textquotedblright{}.

\begin{quote}
Tanti auguri di buon compleanno!
\end{quote}

\begin{quote}
\emph{Happy birthday and all the best!}
\end{quote}

\emph{Fare gli auguri a qualcuno}. Singular \emph{l'augurio}.

\begin{cousinweb}
\emph{auguri} is the plural of \emph{augurio}, a learned borrowing from Latin \emph{augurium} \textquotedblleft{}an omen read by an augur\textquotedblright{} (Lewis \& Short).
\end{cousinweb}

\begin{grammarlens}[title={What you've built}]
One more word for telling, explaining and describing at B1.
\end{grammarlens}

\subsection*{Your turn}
\begin{itemize}
\raggedright
  \item \emph{Say it:} \emph{gli auguri}
  \item \emph{Say it:} \emph{gli auguri}, once more
  \item \emph{Say it:} a sentence of your own with \emph{gli auguri}
  \item \emph{Recall:} say the Italian for the recurring occasion, then the Italian for the anniversary, then say \emph{gli auguri} again
\end{itemize}

\begin{tcolorbox}[breakable,colback=teal!4,colframe=teal!35!black,title={Before you move on}]
What does gli auguri mean? (\textquotedblleft{}Best wishes, greetings\textquotedblright{}.) Say it once more.
\end{tcolorbox}
\section[augurare]{augurare — to wish (someone something)}

\label{lesson:IT-C315-augurare}



\begin{quote}
Before the new one: say the Italian for the anniversary, then the Italian for best wishes.
\end{quote}

\subsection*{You'll want to know: augurare}
\textbf{augurare} — \textquotedblleft{}to wish (someone something)\textquotedblright{}.

\begin{quote}
Ti auguro buon viaggio!
\end{quote}

\begin{quote}
\emph{I wish you a good journey!}
\end{quote}

\emph{Augurare qualcosa a qualcuno}.

\begin{cousinweb}
\emph{augurare} is a learned borrowing from Latin \emph{augurāre} \textquotedblleft{}to read omens\textquotedblright{} (Lewis \& Short), from \emph{augur}, a priest who read the signs.
\end{cousinweb}

\begin{grammarlens}[title={What you've built}]
One more word for telling, explaining and describing at B1.
\end{grammarlens}

\subsection*{Your turn}
\begin{itemize}
\raggedright
  \item \emph{Say it:} \emph{augurare}
  \item \emph{Say it:} \emph{augurare}, once more
  \item \emph{Say it:} a sentence of your own with \emph{augurare}
  \item \emph{Recall:} say the Italian for the anniversary, then the Italian for best wishes, then say \emph{augurare} again
\end{itemize}

\begin{tcolorbox}[breakable,colback=teal!4,colframe=teal!35!black,title={Before you move on}]
What does augurare mean? (\textquotedblleft{}To wish (someone something)\textquotedblright{}.) Say it once more.
\end{tcolorbox}
